\chapter{Materiais e métodos}\label{chap:materiais}

Aqui deve ser descrito com quais elementos e como vai realizar seu trabalho final. Dois itens imprescindíveis são a validação e o cronograma.

\section{Ferramentas}
Descreva aqui as ferramentas a serem utilizadas no desenvolvimento do Projeto. Caso ele utilize algum hardware específico, adicione sua descrição aqui.

\section{Metodologia para o desenvolvimento}
Detalhar cada uma das etapas a serem realizadas no desenvolvimento. Caso faça protótipos, descreva a quantidade deles e a funcionalidade buscada com cada um.

\section{Validação}
Descrever como se pretende validar a solução a ser criada/desenvolvida/melhorada no trabalho. Caso pretenda realizar experimentos, descreva a metodologia de execução e de coleta dos resultados. Descreva como planeja criar a arquitetura para execução dos testes de validação. Deixe claro as métricas a serem usadas na avaliação e validação.

Se a validação contar com formulários, aplicados a terceiros, descrever o(s) formulário(s) a ser utilizado (com referências caso seja validado).

\section{Cronograma}\label{sec:cronograma}

Deve-se incluir no cronograma apenas as atividades a serem desenvolvidas na execução do trabalho final. Não colocar as atividades já desenvolvidas na escrita da proposta.

É sugerido que seja criado um cronograma com indicação de semanas dentro dos meses. Considere que cada mês possui quatro (4) semanas, Utilize um 'X' para delimitar cada um das semanas.

\begin{table}[htb!]
\centering
\caption{Cronograma de atividades da Proposta de Trabalho de Conclusão.}
\label{tab:cronograma}
\begin{tabular}{|l|llllll|}
\hline
\rowcolor[HTML]{C0C0C0} 
  & \multicolumn{6}{c|}{Ano}                                        \\ 
\cline{2-7} 
\rowcolor[HTML]{C0C0C0} 
\multirow{-2}{*}{Atividades} & \multicolumn{1}{c|}{Jul} & \multicolumn{1}{c|}{Ago} & \multicolumn{1}{c|}{Set} & \multicolumn{1}{c|}{Out} & \multicolumn{1}{c|}{Nov} & Dez \\ \hline
Atividade A            & \multicolumn{1}{c|}{XX}   & \multicolumn{1}{c|}{X}    & \multicolumn{1}{c|}{X}   & \multicolumn{1}{c|}{}    & \multicolumn{1}{c|}{}    &   \\ 
\hline
Atividade B            & \multicolumn{1}{c|}{}     & \multicolumn{1}{c|}{X}    & \multicolumn{1}{c|}{}    & \multicolumn{1}{c|}{X}    & \multicolumn{1}{c|}{}    &  \\
\hline
Atividade C            & \multicolumn{1}{c|}{X}    & \multicolumn{1}{c|}{X}    & \multicolumn{1}{c|}{}    & \multicolumn{1}{c|}{}    & \multicolumn{1}{c|}{}    &   \\
\hline
Atividade D            & \multicolumn{1}{c|}{}     & \multicolumn{1}{c|}{XXXX} & \multicolumn{1}{c|}{X}   & \multicolumn{1}{c|}{X}   & \multicolumn{1}{c|}{}    & X \\
\hline
Atividade E            & \multicolumn{1}{c|}{}     & \multicolumn{1}{c|}{}     & \multicolumn{1}{c|}{X}   & \multicolumn{1}{c|}{}    & \multicolumn{1}{c|}{}    &   \\
\hline
Atividade F            & \multicolumn{1}{c|}{}     & \multicolumn{1}{c|}{}     & \multicolumn{1}{c|}{X}   & \multicolumn{1}{c|}{XXX} & \multicolumn{1}{c|}{X}    &   \\
\hline
Atividade G            & \multicolumn{1}{c|}{}     & \multicolumn{1}{c|}{}     & \multicolumn{1}{c|}{}    & \multicolumn{1}{c|}{X}    & \multicolumn{1}{c|}{X}    & X \\
\hline
Elaborar o texto final & \multicolumn{1}{c|}{}     & \multicolumn{1}{c|}{}     & \multicolumn{1}{c|}{}    & \multicolumn{1}{c|}{X}    & \multicolumn{1}{c|}{X}    & X \\
\hline
Banca de defesa        & \multicolumn{1}{c|}{}     & \multicolumn{1}{c|}{}     & \multicolumn{1}{c|}{}    & \multicolumn{1}{c|}{}    & \multicolumn{1}{c|}{}    & X \\
\hline
\end{tabular}
\end{table}

No exemplo da Tabela \autoref{tab:cronograma}, no mês de julho a 'Atividade A' vai consumir duas semanas, no mês de Agosto a 'Atividade D' vai perdurar o mês todo e no mês de Setembro a 'Atividade F' vai levar três semanas.

Procure descrever ao máximo o objetivo, abrangência e os procedimentos de cada atividade, em descrição abaixo da Tabela. \textbf{Uma atividade de revisão bibliográfica NÃO deve constar no cronograma.} Ela faz parte do seu projeto e deve ser realizada ainda durante a preparação da proposta.


